\Askhsh{13056}{4}
\begin{ylh}{2}
    \item 2.1. Οι Πράξεις και οι Ιδιότητές τους
    \item 3.3. Εξισώσεις 2ου Βαθμού
    \item 5.1. Ακολουθίες
\end{ylh}

Οι αριθμοί $1, 3, 6, 10, \dots$ και γενικά αυτοί που είναι δυνατόν, αν παρασταθούν με τελείες, να τοποθετηθούν σε μια τριγωνική διάταξη της μορφής που φαίνεται στον παρακάτω πίνακα λέγονται τριγωνικοί.
\begin{center}
\begin{tikzpicture}[
    dot/.style={circle, fill=maincolor, inner sep=0pt, minimum size=6pt},
    scale=0.8
]
    % T1 = 1
    \begin{scope}[xshift=0cm,anchor=west]
        \node[dot] at (2.5,0) {};
        \node[below] at (2.5,-0.4) {$T_1 = 1$};
    \end{scope}

    % T2 = 3
    \begin{scope}[xshift=2cm,anchor=west]
        \foreach \col in {1,2} {
            \foreach \row in {1,...,\col} {
                \node[dot] at (\col*0.5+1.5, \row*0.5 - 0.5) {};
            }
        }
        \node[below] at (2.5,-0.4) {$T_2 = 3$};
    \end{scope}

    % T3 = 6
    \begin{scope}[xshift=4.5cm,anchor=west]
        \foreach \col in {1,2,3} {
            \foreach \row in {1,...,\col} {
                \node[dot] at (\col*0.5+1, \row*0.5 - 0.5) {};
            }
        }
        \node[below] at (2.5,-0.4) {$T_3 = 6$};
    \end{scope}

    % T4 = 10
    \begin{scope}[xshift=7.5cm,anchor=west]
        \foreach \col in {1,2,3,4} {
            \foreach \row in {1,...,\col} {
                \node[dot] at (\col*0.5+0.5, \row*0.5 - 0.5) {};
            }
        }
        \node[below] at (2.5,-0.4) {$T_4 = 10$};
    \end{scope}

\end{tikzpicture}\captionof{figure}{Άσκηση 13056. Τριγωνικοί αριθμοί}
\end{center}
Αποδεικνύεται ότι ο νιοστός τριγωνικός αριθμός δίνεται από τον τύπο $T_\nu = \dfrac{\nu(\nu+1)}{2}, \nu \in \mathbb{N}^*$.

\begin{alist}
\item Να βρείτε τον 10\textsuperscript{ο} τριγωνικό αριθμό.
\monades{6}
\item Να εξετάσετε αν ο αριθμός $120$ είναι τριγωνικός.
\monades{9}
\item Να αποδείξετε ότι το άθροισμα δυο διαδοχικών τριγωνικών αριθμών είναι ίσο με το τετράγωνο θετικού ακεραίου.
\monades{10}
\end{alist}

